\subsection{3.3. Tìm kiếm chi phí đều (UCS) và hàm chi phí}

\subsubsection{3.3.1. Ý tưởng}

UCS {[}7, §3.4.2{]} luôn mở rộng nút có tổng chi phí \(g(n)\) nhỏ nhất, tính từ trạng thái đầu. Biên là một hàng đợi ưu tiên theo \(g\). BFS mở rộng theo số bước; UCS mở rộng theo chi phí. Vì vậy, khi các bước có chi phí khác nhau, UCS trả về lời giải rẻ nhất, không nhất thiết là lời giải ít bước nhất. Khi mọi bước tốn như nhau, UCS mở rộng theo đúng các lớp độ sâu như BFS và cho cùng số bước (test \texttt{test\_\allowbreak{}ucs\_\allowbreak{}uniform\_\allowbreak{}cost\_\allowbreak{}matches\_\allowbreak{}par\_\allowbreak{}on\_\allowbreak{}all\_\allowbreak{}33\_\allowbreak{}levels}).

Đề yêu cầu báo cáo định nghĩa và biện minh hàm chi phí, và lưu ý rằng nếu mọi bước tốn 1 thì UCS trùng với BFS {[}1, mục 3.4{]}. Vì vậy nhóm dùng một hàm chi phí không đồng nhất. Mục 3.3.2 định nghĩa hàm chi phí của nhóm, mục 3.3.3 giải thích lý do của từng mức. Trạng thái, hành động, hàm chuyển và kiểm tra đích dùng chung với ba thuật toán còn lại (chương 2).

\subsubsection{3.3.2. Hàm chi phí: định nghĩa}

Chi phí của bước \(s \xrightarrow{a} s'\) gồm hai phần: phần phụ thuộc tư thế trước và sau khi lăn, và phụ phí khi khối đè lên ô dễ vỡ:

\[
c(s, a, s') = c_{\text{lăn}}\big(o(s), o(s')\big) + \begin{cases} 0{,}5 & \text{nếu có ô } x \in \text{cells}(s') \text{ là ô dễ vỡ,} \\ 0 & \text{nếu không,} \end{cases}
\]

trong đó \(o(s)\) là tư thế của khối ở trạng thái \(s\), và \(\text{cells}(s')\) là các ô khối chiếm sau bước đi. Bảng 3.3.1 cho \(c_{\text{lăn}}\).

\begin{longtable}[]{@{}
  >{\raggedright\arraybackslash}p{(\linewidth - 8\tabcolsep) * \real{0.2850}}
  >{\raggedright\arraybackslash}p{(\linewidth - 8\tabcolsep) * \real{0.2705}}
  >{\raggedright\arraybackslash}p{(\linewidth - 8\tabcolsep) * \real{0.1111}}
  >{\raggedright\arraybackslash}p{(\linewidth - 8\tabcolsep) * \real{0.1111}}
  >{\raggedright\arraybackslash}p{(\linewidth - 8\tabcolsep) * \real{0.2222}}@{}}
\caption{\textbf{Bảng 3.3.1.} Chi phí lăn theo tư thế trước và sau bước đi. \(\Delta h\) là độ cao trọng tâm phải nâng thêm để khối lật qua cạnh; \(m\) là khối lượng, lấy khối 1×1×2 có \(m = 1\).}\tabularnewline
\toprule\noalign{}
\begin{minipage}[b]{\linewidth}\raggedright
\textbf{Bước (ví dụ)}
\end{minipage} & \begin{minipage}[b]{\linewidth}\raggedright
\textbf{Trọng tâm: lúc đầu → đỉnh cung lật}
\end{minipage} & \begin{minipage}[b]{\linewidth}\raggedright
\textbf{\(\Delta h\)}
\end{minipage} & \begin{minipage}[b]{\linewidth}\raggedright
\textbf{Công nâng \(m \cdot \Delta h\)}
\end{minipage} & \begin{minipage}[b]{\linewidth}\raggedright
\textbf{\(c_{\text{lăn}}\)}
\end{minipage} \\
\midrule\noalign{}
\endfirsthead
\toprule\noalign{}
\begin{minipage}[b]{\linewidth}\raggedright
\textbf{Bước (ví dụ)}
\end{minipage} & \begin{minipage}[b]{\linewidth}\raggedright
\textbf{Trọng tâm: lúc đầu → đỉnh cung lật}
\end{minipage} & \begin{minipage}[b]{\linewidth}\raggedright
\textbf{\(\Delta h\)}
\end{minipage} & \begin{minipage}[b]{\linewidth}\raggedright
\textbf{Công nâng \(m \cdot \Delta h\)}
\end{minipage} & \begin{minipage}[b]{\linewidth}\raggedright
\textbf{\(c_{\text{lăn}}\)}
\end{minipage} \\
\midrule\noalign{}
\endhead
\bottomrule\noalign{}
\endlastfoot
Đứng → nằm (khối đứng lăn R) & \(1 \to \sqrt{5}/2 \approx 1{,}118\) & 0,118 & 0,118 & 1,12 \\
Nằm → đứng (khối nằm ngang lăn R) & \(0{,}5 \to \sqrt{5}/2 \approx 1{,}118\) & 0,618 & 0,618 & 1,62 \\
Nằm → nằm, lăn quanh trục dài (khối nằm ngang lăn U) & \(0{,}5 \to \sqrt{2}/2 \approx 0{,}707\) & 0,207 & 0,207 & 1,21 \\
Khối con 1×1×1 đi một ô, \(m = 1/2\) (lăn U1) & \(0{,}5 \to \sqrt{2}/2 \approx 0{,}707\) & 0,207 & 0,104 & 1,10 \\
Khối con đi rồi hai khối con ghép thành khối nằm & như khối con đi một ô & 0,207 & 0,104 & 1,10 \\
Nằm → đứng lên công tắc tách, rồi tách & như nằm → đứng & 0,618 & 0,618 & 1,62 \\
\end{longtable}

Công thức chung là \(c_{\text{lăn}} = 1 + m \cdot \Delta h\), làm tròn hai chữ số thập phân. Hàm cài đặt là \texttt{compute\_\allowbreak{}step\_\allowbreak{}cost} trong \texttt{Source/\allowbreak{}bloxorz/\allowbreak{}costs.\allowbreak{}py}; \texttt{Problem.\allowbreak{}step\_\allowbreak{}cost(s, a, s')} gọi hàm này. Bảng 2.7 ở mục 2.8 tóm tắt cùng các mức này.

Trên 340 414 cạnh tới được của 33 màn gốc, chi phí một bước nằm trong khoảng 1,10 (khối con đi) đến 1,71 (khối nằm lăn quanh trục dài lên ô dễ vỡ: 1,21 + 0,5). Khối nguyên đứng trên ô dễ vỡ thì rơi (mục 2.1), nên không có bước ``nằm → đứng'' nào chịu phụ phí ô dễ vỡ.

\subsubsection{3.3.3. Lý do chọn hàm chi phí}

\textbf{Phần cơ sở 1 cho mỗi bước.} Mỗi lần lăn là một thao tác của người chơi và tốn một khoảng thời gian như nhau. Mức cơ sở này cũng bảo đảm mọi bước tốn ít nhất 1. A* cần điều kiện đó: hai heuristic ở mục 3.4 đếm cận dưới của \textbf{số bước} còn lại, nên chúng chỉ chấp nhận được khi mỗi bước tốn ít nhất 1 (chứng minh ở mục 3.4.4).

\textbf{Phần công nâng trọng tâm.} Khối lăn bằng cách lật quanh một cạnh đáy. Trong lúc lật, trọng tâm đi trên một cung tròn tâm là cạnh đó, và lên cao nhất khi nằm thẳng trên cạnh. Bán kính cung bằng khoảng cách từ cạnh tới trọng tâm:

\begin{itemize}
\tightlist
\item
  Khối đứng lật quanh cạnh đáy dài 1: trọng tâm ở độ cao 1, cách cạnh \(\sqrt{0{,}5^2 + 1^2} = \sqrt{5}/2\).
\item
  Khối nằm dựng lên quanh cạnh ngắn ở đầu khối: trọng tâm ở độ cao 0,5, cách cạnh \(\sqrt{1^2 + 0{,}5^2} = \sqrt{5}/2\).
\item
  Khối nằm lăn quanh cạnh dài, hoặc khối con 1×1×1 lăn: trọng tâm ở độ cao 0,5, cách cạnh \(\sqrt{0{,}5^2 + 0{,}5^2} = \sqrt{2}/2\).
\end{itemize}

Công cần để lật khối là công nâng trọng tâm từ độ cao lúc đầu lên đỉnh cung, \(m \cdot g \cdot \Delta h\). Phần trọng tâm hạ xuống sau đó không được tính, vì khối tự đổ xuống. Lấy \(m \cdot g = 1\) cho khối nguyên, ta có cột ``Công nâng'' ở Bảng 3.3.1. Khối con chỉ nặng bằng nửa khối nguyên, nên công của nó bằng một nửa. Thứ tự các mức khớp với trực giác khi chơi: dựng khối từ nằm lên đứng là động tác nặng nhất (1,62), lật khối đứng cho nằm xuống thì nhẹ (1,12).

\textbf{Bước tách khối và bước ghép khối.} Khối chỉ tách khi đứng lên công tắc tách, tức là bước đó là một bước ``nằm → đứng'' (Bảng 2.4); việc tách diễn ra sau khi khối đã đứng trên công tắc. Vì vậy bước tách tốn như ``nằm → đứng'', 1,62. Bước ghép là một bước khối con đi một ô; hai khối con tự ghép khi chạm nhau, không phải nâng thêm gì, nên tốn 1,10. Không có hai mức này thì bước tách rơi về mức mặc định 1,0 trong code, rẻ hơn mọi bước khác, và UCS sẽ ưu tiên đi qua công tắc tách một cách vô lý. Test \texttt{test\_\allowbreak{}every\_\allowbreak{}reachable\_\allowbreak{}edge\_\allowbreak{}uses\_\allowbreak{}energy\_\allowbreak{}table} duyệt mọi cạnh tới được của 33 màn và bảo đảm không cạnh nào rơi về mức mặc định.

\textbf{Phụ phí ô dễ vỡ 0,5.} Khối được phép nằm trên ô dễ vỡ, nhưng chỉ cần một bước dựng khối sai là ô vỡ và thua. Phụ phí 0,5 biểu diễn rủi ro này: khi có nhiều đường cùng số bước, UCS chọn đường ít đè lên ô dễ vỡ hơn. Mức 0,5 bằng một nửa phần cơ sở, nên để tránh được một lần đè ô dễ vỡ, UCS chưa chấp nhận đi thêm cả một bước (tốn ít nhất 1,10). Phụ phí vì vậy chỉ phân định giữa các đường gần bằng nhau, không kéo UCS đi đường vòng xa. Đây cũng là ví dụ mà đề gợi ý: đi lên ô dễ vỡ thì ``đắt'' vì rủi ro {[}1, mục 3.4{]}.

\textbf{Vì sao không giảm chi phí khi bấm công tắc.} Đề cũng gợi ý có thể cho việc bấm công tắc ``rẻ''. Nhóm không làm vậy: một bước bấm công tắc rẻ hơn 1 sẽ vi phạm điều kiện mọi bước tốn ít nhất 1, và hai heuristic đếm số bước của A* không còn chấp nhận được (mục 3.4.4). Bấm công tắc vẫn là một lần lăn, nên tính theo cùng Bảng 3.3.1.

\textbf{Các tính chất dùng ở những phần khác.}

\begin{enumerate}
\def\labelenumi{\arabic{enumi}.}
\tightlist
\item
  Mọi bước tốn ít nhất 1,10 \textgreater{} 1: điều kiện của chứng minh ở mục 3.4.4 thỏa, và không có chu trình chi phí 0, nên UCS dừng và tối ưu (mục 3.3.5).
\item
  Chi phí chỉ phụ thuộc tư thế trước, tư thế sau và các ô dưới khối sau bước đi, nên tính được trong \(O(1)\) và không cần thêm gì vào trạng thái.
\item
  Cả năm cấu hình ở chương 5 tính tổng chi phí lời giải bằng cùng hàm này, kể cả BFS và DFS vốn không dùng chi phí khi tìm. Nhờ vậy cột tổng chi phí so sánh được giữa các thuật toán.
\end{enumerate}

\textbf{Hàm chi phí làm UCS chọn đường khác BFS ở đâu.} Trên 33 màn gốc, UCS và BFS luôn có cùng số bước, nhưng có tổng chi phí khác nhau ở 2 màn (Bảng 3.3.2). Ở màn 17, đường của UCS dựng khối đứng ít hơn một lần, thay bằng hai lần lăn khối nằm: bớt một bước ``nằm → đứng'' (1,62) và một bước ``đứng → nằm'' (1,12), thêm hai bước ``nằm → nằm'' (2 × 1,21), tức rẻ hơn 0,32. Ở màn 23, đường của UCS đè lên ô dễ vỡ 21 lần thay vì 23 lần, nên rẻ hơn 2 × 0,5 = 1,00. Test \texttt{test\_\allowbreak{}branching\_\allowbreak{}decisions\_\allowbreak{}on\_\allowbreak{}levels\_\allowbreak{}17\_\allowbreak{}and\_\allowbreak{}23} kiểm tra hai trường hợp này.

\begin{longtable}[]{@{}
  >{\raggedright\arraybackslash}p{(\linewidth - 8\tabcolsep) * \real{0.0694}}
  >{\raggedright\arraybackslash}p{(\linewidth - 8\tabcolsep) * \real{0.0833}}
  >{\raggedright\arraybackslash}p{(\linewidth - 8\tabcolsep) * \real{0.1111}}
  >{\raggedright\arraybackslash}p{(\linewidth - 8\tabcolsep) * \real{0.1111}}
  >{\raggedright\arraybackslash}p{(\linewidth - 8\tabcolsep) * \real{0.6250}}@{}}
\caption{\textbf{Bảng 3.3.2.} Hai màn có lời giải BFS khác lời giải UCS về tổng chi phí (cùng số bước).}\tabularnewline
\toprule\noalign{}
\begin{minipage}[b]{\linewidth}\raggedright
\textbf{Màn}
\end{minipage} & \begin{minipage}[b]{\linewidth}\raggedright
\textbf{Số bước}
\end{minipage} & \begin{minipage}[b]{\linewidth}\raggedright
\textbf{Tổng chi phí BFS}
\end{minipage} & \begin{minipage}[b]{\linewidth}\raggedright
\textbf{Tổng chi phí UCS}
\end{minipage} & \begin{minipage}[b]{\linewidth}\raggedright
\textbf{Khác nhau ở đâu}
\end{minipage} \\
\midrule\noalign{}
\endfirsthead
\toprule\noalign{}
\begin{minipage}[b]{\linewidth}\raggedright
\textbf{Màn}
\end{minipage} & \begin{minipage}[b]{\linewidth}\raggedright
\textbf{Số bước}
\end{minipage} & \begin{minipage}[b]{\linewidth}\raggedright
\textbf{Tổng chi phí BFS}
\end{minipage} & \begin{minipage}[b]{\linewidth}\raggedright
\textbf{Tổng chi phí UCS}
\end{minipage} & \begin{minipage}[b]{\linewidth}\raggedright
\textbf{Khác nhau ở đâu}
\end{minipage} \\
\midrule\noalign{}
\endhead
\bottomrule\noalign{}
\endlastfoot
17 & 106 & 140,42 & 140,10 & UCS: 9 lần dựng khối đứng (BFS: 10), 25 lần lăn khối nằm ngang (BFS: 23) \\
23 & 75 & 101,58 & 100,58 & UCS: 21 bước đè ô dễ vỡ (BFS: 23); số lần đổi tư thế như nhau \\
\end{longtable}

\subsubsection{3.3.4. Mã giả}

\begin{verbatim}
UCS(problem):
    s0 ← problem.initial_state()
    frontier ← hàng đợi ưu tiên [(g = 0, thứ_tự = 0, s0)];  best_g[s0] ← 0
    lặp khi frontier khác rỗng:
        (g, _, s) ← lấy ra mục có (g, thứ_tự) nhỏ nhất    ▷ g bằng nhau: vào trước ra trước
        nếu g > best_g[s]: bỏ qua                         ▷ mục cũ (lazy deletion)
        nếu problem.is_goal(s): trả về dựng_đường(cha, s), g   ▷ kiểm tra đích khi LẤY RA
        expanded ← expanded + 1
        với mỗi a trong problem.actions(s):               ▷ thứ tự U, D, L, R (đã tách: U1 … R1, U2 … R2)
            s' ← problem.result(s, a)
            nếu s' = None: bỏ qua                         ▷ khối rơi, không phải nút con
            generated ← generated + 1
            g' ← g + problem.step_cost(s, a, s')
            nếu g' ≥ best_g[s'] (nếu đã có): bỏ qua         ▷ không rẻ hơn đường đã biết
            best_g[s'] ← g';  cha[s'] ← (s, a)
            đưa (g', thứ_tự mới, s') vào frontier
            max_frontier ← max(max_frontier, |frontier|)
    trả về NO_SOLUTION
\end{verbatim}

Mã nguồn: \texttt{Source/\allowbreak{}solvers/\allowbreak{}ucs.\allowbreak{}py} (lớp \texttt{UCS}, hàm \texttt{\_\allowbreak{}path}) và \texttt{Source/\allowbreak{}bloxorz/\allowbreak{}costs.\allowbreak{}py}.

\subsubsection{3.3.5. Tính đầy đủ, tối ưu và độ phức tạp}

Gọi \(S\) là tập trạng thái tới được, \(b \le 8\) là số hành động tối đa ở một trạng thái, \(C^*\) là tổng chi phí của lời giải tối ưu và \(\varepsilon = 1{,}10\) là chi phí bước nhỏ nhất.

\begin{itemize}
\tightlist
\item
  \textbf{Đầy đủ.} \(S\) hữu hạn (mục 2.3) và mọi bước tốn ít nhất \(\varepsilon > 0\). Một trạng thái chỉ được đưa lại vào biên khi có đường rẻ hơn hẳn, và mục lấy ra có \(g\) lớn hơn \texttt{best\_\allowbreak{}g} của trạng thái đó thì bị bỏ qua. Với chi phí dương, mỗi trạng thái được mở rộng nhiều nhất một lần: lần đầu lấy ra với \(g\) bằng \texttt{best\_\allowbreak{}g} thì \(g\) đó đã nhỏ nhất. Vì vậy vòng lặp dừng: có lời giải thì tìm ra, không có thì biên rỗng và trả \texttt{no\_\allowbreak{}solution}.
\item
  \textbf{Tối ưu theo tổng chi phí.} UCS lấy các nút ra theo thứ tự \(g\) không giảm. Khi trạng thái đích được lấy ra với chi phí \(g\), mọi đường khác tới đích đều phải đi qua một nút còn trong biên, có chi phí ít nhất \(g\); chi phí bước dương nên đi tiếp chỉ làm đường đắt thêm. Vì vậy \(g = C^*\) {[}7, §3.4.2{]}. Điều này chỉ đúng khi kiểm tra đích lúc \textbf{lấy ra}; nếu kiểm tra khi sinh nút như BFS, UCS có thể trả về một đường đắt hơn vừa được sinh trước. Trên 33 màn, tổng chi phí của UCS bằng của A* ở 33/33 màn, và không lớn hơn của BFS, DFS ở màn nào (chương 5).
\item
  \textbf{Thời gian.} Theo lý thuyết, UCS tốn \(O(b^{1 + \lfloor C^*/\varepsilon \rfloor})\) {[}7, §3.4.2{]}. Với tập \texttt{best\_\allowbreak{}g}, số lần mở rộng còn bị chặn bởi \(|S|\), số lần đưa vào heap bởi \(b \cdot |S|\), mỗi thao tác heap tốn \(O(\log(b|S|))\), nên tổng thời gian là \(O(b|S| \log(b|S|))\).
\item
  \textbf{Bộ nhớ.} \(O(b \cdot |S|)\) cho heap (tính cả mục cũ chờ bỏ qua) và \(O(|S|)\) cho \texttt{best\_\allowbreak{}g} và bảng cha.
\end{itemize}

\subsubsection{3.3.6. Các điểm riêng khi cài đặt}

\begin{enumerate}
\def\labelenumi{\arabic{enumi}.}
\tightlist
\item
  \textbf{Hàng đợi ưu tiên \texttt{heapq}, mục là \texttt{(g, thứ\_\allowbreak{}tự, s)}.} \texttt{thứ\_\allowbreak{}tự} lấy từ \texttt{itertools.\allowbreak{}count()}, tăng dần. Khi hai mục có cùng \(g\), Python so \texttt{thứ\_\allowbreak{}tự} mà không phải so hai trạng thái (lớp \texttt{State} không định nghĩa phép so sánh thứ tự), và mục vào trước được lấy ra trước. Nhờ vậy kết quả tất định.
\item
  \textbf{Lazy deletion thay cho giảm khóa.} \texttt{heapq} không có thao tác giảm khóa. Khi tìm được đường rẻ hơn tới một trạng thái đã ở trong biên, UCS đưa thêm một mục mới và cập nhật \texttt{best\_\allowbreak{}g}; mục cũ vẫn nằm trong heap và bị bỏ qua khi lấy ra. Mục bị bỏ qua không tính vào \texttt{expanded}; \texttt{max\_\allowbreak{}frontier} tính cả các mục cũ.
\item
  \textbf{Kiểm tra đích khi lấy ra.} Đây là khác biệt cố ý với BFS và DFS (mục 3.1.3): kiểm tra khi lấy ra mới giữ được tính tối ưu theo chi phí. Cái giá là UCS phải mở rộng mọi trạng thái có \(g < C^*\) trước khi dừng (mục 5.4.3).
\item
  \textbf{Chi phí lấy qua \texttt{Problem.\allowbreak{}step\_\allowbreak{}cost}.} UCS không biết công thức chi phí. Khi tạo \texttt{Problem} không truyền hàm chi phí thì mỗi bước tốn 1, và UCS trả về cùng số bước với BFS; khi truyền \texttt{compute\_\allowbreak{}step\_\allowbreak{}cost} thì dùng Bảng 3.3.1. Giao diện và script thực nghiệm đều truyền \texttt{compute\_\allowbreak{}step\_\allowbreak{}cost}.
\item
  \textbf{So sánh chi phí có làm tròn.} Mỗi bước được làm tròn hai chữ số, nhưng tổng là số thực dấu phẩy động. Các test và công cụ kiểm tra so tổng chi phí theo sai lệch tương đối \(10^{-9}\), không so bằng tuyệt đối.
\item
  \textbf{Đếm chỉ số theo quy ước chung} (mục 3.1.3): \texttt{generated} đếm mọi con hợp lệ trước khi xét \texttt{best\_\allowbreak{}g}; \texttt{visited} là số khóa của \texttt{best\_\allowbreak{}g}, tính cả trạng thái đầu.
\end{enumerate}

\begin{longtable}[]{@{}
  >{\raggedright\arraybackslash}p{(\linewidth - 16\tabcolsep) * \real{0.0880}}
  >{\raggedright\arraybackslash}p{(\linewidth - 16\tabcolsep) * \real{0.1019}}
  >{\raggedright\arraybackslash}p{(\linewidth - 16\tabcolsep) * \real{0.1019}}
  >{\raggedright\arraybackslash}p{(\linewidth - 16\tabcolsep) * \real{0.1296}}
  >{\raggedright\arraybackslash}p{(\linewidth - 16\tabcolsep) * \real{0.1157}}
  >{\raggedright\arraybackslash}p{(\linewidth - 16\tabcolsep) * \real{0.1157}}
  >{\raggedright\arraybackslash}p{(\linewidth - 16\tabcolsep) * \real{0.1296}}
  >{\raggedright\arraybackslash}p{(\linewidth - 16\tabcolsep) * \real{0.1019}}
  >{\raggedright\arraybackslash}p{(\linewidth - 16\tabcolsep) * \real{0.1157}}@{}}
\caption{\textbf{Bảng 3.3.3.} UCS so với BFS trên một số màn gốc (hàm chi phí ở Bảng 3.3.1). Số đếm tất định, không phụ thuộc máy chạy.}\tabularnewline
\toprule\noalign{}
\begin{minipage}[b]{\linewidth}\raggedright
\textbf{Màn}
\end{minipage} & \begin{minipage}[b]{\linewidth}\raggedright
\textbf{Nhóm}
\end{minipage} & \begin{minipage}[b]{\linewidth}\raggedright
\textbf{Số bước}
\end{minipage} & \begin{minipage}[b]{\linewidth}\raggedright
\textbf{Tổng chi phí}
\end{minipage} & \begin{minipage}[b]{\linewidth}\raggedright
\textbf{Nút mở rộng UCS}
\end{minipage} & \begin{minipage}[b]{\linewidth}\raggedright
\textbf{Nút mở rộng BFS}
\end{minipage} & \begin{minipage}[b]{\linewidth}\raggedright
\textbf{Nút sinh UCS}
\end{minipage} & \begin{minipage}[b]{\linewidth}\raggedright
\textbf{Biên lớn nhất UCS}
\end{minipage} & \begin{minipage}[b]{\linewidth}\raggedright
\textbf{Đã thăm UCS}
\end{minipage} \\
\midrule\noalign{}
\endfirsthead
\toprule\noalign{}
\begin{minipage}[b]{\linewidth}\raggedright
\textbf{Màn}
\end{minipage} & \begin{minipage}[b]{\linewidth}\raggedright
\textbf{Nhóm}
\end{minipage} & \begin{minipage}[b]{\linewidth}\raggedright
\textbf{Số bước}
\end{minipage} & \begin{minipage}[b]{\linewidth}\raggedright
\textbf{Tổng chi phí}
\end{minipage} & \begin{minipage}[b]{\linewidth}\raggedright
\textbf{Nút mở rộng UCS}
\end{minipage} & \begin{minipage}[b]{\linewidth}\raggedright
\textbf{Nút mở rộng BFS}
\end{minipage} & \begin{minipage}[b]{\linewidth}\raggedright
\textbf{Nút sinh UCS}
\end{minipage} & \begin{minipage}[b]{\linewidth}\raggedright
\textbf{Biên lớn nhất UCS}
\end{minipage} & \begin{minipage}[b]{\linewidth}\raggedright
\textbf{Đã thăm UCS}
\end{minipage} \\
\midrule\noalign{}
\endhead
\bottomrule\noalign{}
\endlastfoot
1 & G1 & 7 & 9,11 & 66 & 48 & 183 & 24 & 87 \\
8 & G5 & 10 & 12,60 & 594 & 412 & 3 343 & 135 & 656 \\
17 & G3 & 106 & 140,10 & 988 & 977 & 2 429 & 47 & 992 \\
20 & G5 & 56 & 70,68 & 4 438 & 3 983 & 23 550 & 366 & 4 800 \\
23 & G5 & 75 & 100,58 & 3 755 & 3 692 & 20 060 & 190 & 3 831 \\
29 & G3 & 104 & 140,24 & 1 811 & 1 809 & 4 007 & 48 & 1 816 \\
32 & G3 & 129 & 168,57 & 626 & 626 & 1 308 & 18 & 632 \\
\end{longtable}

UCS mở rộng nhiều nút hơn BFS ở các màn có lời giải ngắn so với kích thước không gian (màn 1, 8): BFS dừng ngay khi sinh ra đích, còn UCS phải mở rộng hết các trạng thái rẻ hơn lời giải. Ở các màn có lời giải dài và hành lang hẹp (màn 29, 32), gần như mọi trạng thái đều rẻ hơn lời giải, nên hai thuật toán mở rộng gần bằng nhau. Nhận xét đầy đủ ở mục 5.4.3.

\subsubsection{3.3.7. Kiểm thử}

\texttt{Source/\allowbreak{}tests/\allowbreak{}test\_\allowbreak{}ucs.\allowbreak{}py}, \texttt{Source/\allowbreak{}tests/\allowbreak{}test\_\allowbreak{}costs.\allowbreak{}py} và \texttt{Source/\allowbreak{}tests/\allowbreak{}test\_\allowbreak{}costs\_\allowbreak{}split.\allowbreak{}py} gồm:

\begin{itemize}
\tightlist
\item
  \texttt{test\_\allowbreak{}start\_\allowbreak{}is\_\allowbreak{}goal\_\allowbreak{}returns\_\allowbreak{}empty\_\allowbreak{}path}, \texttt{test\_\allowbreak{}counts\_\allowbreak{}on\_\allowbreak{}corridor\_\allowbreak{}by\_\allowbreak{}hand}, \texttt{test\_\allowbreak{}no\_\allowbreak{}solution\_\allowbreak{}when\_\allowbreak{}gap\_\allowbreak{}is\_\allowbreak{}two\_\allowbreak{}tiles}: trạng thái đầu là đích, đếm chỉ số trên hành lang so với kết quả tính tay, màn không có lời giải;
\item
  \texttt{test\_\allowbreak{}all\_\allowbreak{}mazes\_\allowbreak{}match\_\allowbreak{}expected\_\allowbreak{}cost}: bốn mê cung mẫu (có mê cung với ô chi phí khác 1 và mê cung không có lời giải), so với chi phí tối ưu biết trước;
\item
  \texttt{test\_\allowbreak{}ucs\_\allowbreak{}uniform\_\allowbreak{}cost\_\allowbreak{}matches\_\allowbreak{}par\_\allowbreak{}on\_\allowbreak{}all\_\allowbreak{}33\_\allowbreak{}levels}: chi phí bằng 1 thì UCS ra đúng số bước tối ưu ở 33 màn, như BFS;
\item
  \texttt{test\_\allowbreak{}ucs\_\allowbreak{}energy\_\allowbreak{}cost\_\allowbreak{}optimality\_\allowbreak{}on\_\allowbreak{}all\_\allowbreak{}33\_\allowbreak{}levels}: với hàm chi phí ở Bảng 3.3.1, tổng chi phí của UCS không lớn hơn của BFS ở mọi màn, và lời giải phát lại trên engine tới đích;
\item
  \texttt{test\_\allowbreak{}branching\_\allowbreak{}decisions\_\allowbreak{}on\_\allowbreak{}levels\_\allowbreak{}17\_\allowbreak{}and\_\allowbreak{}23}: hai màn UCS chọn đường rẻ hơn BFS (Bảng 3.3.2);
\item
  \texttt{test\_\allowbreak{}all\_\allowbreak{}costs\_\allowbreak{}at\_\allowbreak{}least\_\allowbreak{}one} và các test từng mức (\texttt{test\_\allowbreak{}standing\_\allowbreak{}to\_\allowbreak{}lying\_\allowbreak{}cost}, \texttt{test\_\allowbreak{}lying\_\allowbreak{}to\_\allowbreak{}standing\_\allowbreak{}cost}, \texttt{test\_\allowbreak{}lying\_\allowbreak{}to\_\allowbreak{}lying\_\allowbreak{}cost}, \texttt{test\_\allowbreak{}split\_\allowbreak{}cubes\_\allowbreak{}cost}, \texttt{test\_\allowbreak{}fragile\_\allowbreak{}tile\_\allowbreak{}penalty});
\item
  \texttt{test\_\allowbreak{}lying\_\allowbreak{}to\_\allowbreak{}split\_\allowbreak{}costs\_\allowbreak{}like\_\allowbreak{}lying\_\allowbreak{}to\_\allowbreak{}standing}, \texttt{test\_\allowbreak{}every\_\allowbreak{}reachable\_\allowbreak{}edge\_\allowbreak{}uses\_\allowbreak{}energy\_\allowbreak{}table}, \texttt{test\_\allowbreak{}split\_\allowbreak{}step\_\allowbreak{}and\_\allowbreak{}path\_\allowbreak{}total\_\allowbreak{}on\_\allowbreak{}level\_\allowbreak{}08}: bước tách tốn 1,62, không cạnh tới được nào rơi về mức mặc định, và tổng chi phí của lời giải màn 8 khớp khi phát lại.
\end{itemize}
